\section{Limitations and Threats to Validity}
\label{sec:limitations}

\paragraph{No mechanized proof.}
The propositions and contracts are encoded as executable properties and checked over finite
generators and artifacts.  A passing FsCheck property is not a proof over all
behaviors, schedulers, or logs.  Future work may mechanize selected definitions
after the operational contracts stabilize.

\paragraph{No minimal confluence condition.}
We falsify general confluence and declared-write disjointness.  We check a
restricted isolated-writer family.  We do not establish a necessary or
sufficient general condition over reads, writes, emissions, triggers, and state
domains.  Calling the current restricted condition ``the'' restoring condition
would overstate the result.

\paragraph{Limited scheduler space.}
The executable suite uses four deterministic scheduler extremes.  They expose
the preserved counterexamples but do not enumerate every interleaving for an
arbitrary behavior set.  Relation ordering is discussed from source audit but
not modeled as a separate production activation dimension.  Seeded randomized
and bounded exhaustive interleaving exploration remain future checks.

\paragraph{Adapter assumptions.}
Real-log classification relies on explicit event names and payload fields.  The
adapter is deliberately conservative, but mistakes remain possible.  Model
calls are treated as OneShot for cost/oracle semantics; another property might
classify a read-only deterministic local model differently.  A derived content
hash proves captured bytes, not provider authenticity, commitment, or
idempotence.

The v0 effect descriptor has no per-effect observable-set label.  Its
CounterfactualWorld result therefore uses the selected profile's entire
recognized external surface as $\Omega_{profile}$.  Narrower claims require
adapter-supplied observable tags and new conformance vectors.

Unknown footprint fails closed only after the adapter recognizes an event as an
effect.  Unclassified source events remain reported signals.  The all-cut Sound
counts therefore depend on an audited profile declaration that those event
types do not hide external interactions.  A newly introduced effect type can
make an old verdict false until the adapter is updated.

\paragraph{Partial failure and compensation scope.}
Failed remote operations may have partially committed before reporting an
error.  The assessor conservatively conditions or blocks discarded failures by
footprint, but it cannot inspect the remote system.  Likewise, successful
compensation does not erase a temporal exposure window or third-party secondary
effects.  A caller whose $\Omega$ includes those histories needs stronger
reconciliation evidence than the current effect descriptor records.

\paragraph{Whole-log malformedness policy.}
The current assessor treats structural malformedness anywhere in the source log
as a blocker for every cut, including projection.  This is fail-closed and
appropriate for the evidence study, but it is stricter than a theorem about an
independently authenticated well-formed prefix.  The projection suffix identity is
therefore stated only for well-formed ordered logs.

\paragraph{Synthetic content.}
The Synthetic Players capsule contains archived experimental runs, but the
present analysis concerns runtime traces rather than human ground truth.  The
bridge study uses deterministic offline mocks and explicitly makes no real-model
performance claim.  Its agreement numbers illustrate equivalence relations,
not predictive validity.

\paragraph{Fork coverage.}
The 121 observed forks are useful positive evidence, but all cut at
\code{round.played} and retain no classified external request.  No observed
fork in the public capsule crosses a recorded oracle or irreversible effect
boundary.  The difficult cases are exercised by all-cut analysis and fixtures,
not by an actual public fork with a no-reexecution receipt.  Moreover, all
160,076 Sound public ExternalContinuation cuts precede the first classified
effect.  The current target-environment premise is caller-supplied and has no
implemented attestation-discharge path, so neither observed forks nor those
Sound counts validate one.

\paragraph{Bridge provenance.}
The bridge artifact is locally content-addressed and now bound to local source
and release-bundle commits with a standalone verifier.  The repository still
has no public remote, so a fresh external reviewer cannot fetch those revisions.
The result remains illustrative and excluded from the abstract until the
bundle is published.  This limitation is particularly important because
provenance is part of the paper's thesis.

\paragraph{Production behavior inventory.}
The ActiveGraph source establishes FIFO and registration-order scheduling, but
we did not recover a complete deployed behavior dependency inventory.  We
therefore cannot say whether production behaviors satisfy a future
non-interference condition.  Static source inspection of the scheduler is not
production confluence validation.

\paragraph{No external runtime corpus.}
The public corpus, observed forks, bridge implementation, and local executive
study all descend from ActiveGraph.  The language-neutral vectors support other
implementations, but no LangGraph, Temporal, or other independent runtime trace
has yet been normalized.  Generalization beyond the source family is therefore
conceptual rather than empirically demonstrated.

\paragraph{No authorization or information-flow result.}
Effect descriptors classify replay and environmental conditions; they do not
grant capabilities or prove that a caller is authorized.  The artifact also
does not track secrecy labels or prove noninterference.  Those topics are
covered by distinct systems and calculi and remain explicit non-goals.

\paragraph{No performance benchmark.}
The all-cut implementation is efficient enough for the reported corpora, but
we make no throughput, storage, or latency claim.  ``Cheap'' in this paper is
analyzed semantically.  A separate benchmark would be needed for production
performance.
